
\documentclass[11pt]{article}

\usepackage[margin=1in]{geometry}
\usepackage[T1]{fontenc}
\usepackage[utf8]{inputenc}
\usepackage{lmodern}
\usepackage{microtype}
\usepackage{booktabs}
\usepackage{longtable}
\usepackage{array}
\usepackage{tabularx}
\usepackage{enumitem}
\usepackage{hyperref}
\usepackage{xcolor}
\usepackage{amsmath}
\usepackage{amssymb}

\hypersetup{
  colorlinks=true,
  linkcolor=blue!50!black,
  citecolor=blue!50!black,
  urlcolor=blue!50!black
}

\title{\textbf{Latent Recursion in Small Language Models:}\\
A Controlled Pilot Study of Four Inference-Time Hallucination Engines}
\author{Woflfren}
\date{April 2026}

\begin{document}
\maketitle

\begin{abstract}
This paper rebuilds the ``hallucination engines / latent recursion'' project on an evidence-first footing. We report a controlled pilot study based only on a single executed suite archived on 2026-04-03, containing four text-only inference-time recursion systems: Recursive Token Self-Mirroring (RTSM), Activation Gradient Climbing (AGC), Latent Vector Transform Composition (LVTC), and Self-Reflective Prompt Feedback (SRPF). The suite used three small Hugging Face models and one common seed prompt concerning an underwater machine whose message becomes increasingly strange and self-aware. 

All four runs completed and produced structured JSON/TXT outputs. The resulting behaviors were sharply different. RTSM preserved the original motif most strongly, but suffered catastrophic late-stage repetition and prompt leakage, with type-token ratio falling to 0.510 and repeated bigrams rising to 235 at step 4. AGC remained lexically diverse but drifted off-topic almost immediately, with mean adjacent-output Jaccard similarity of only 0.013 and zero final seed-term retention. LVTC produced moderately diverse text but repeatedly slipped into instruction-following/meta-task discourse, ending with four heuristic meta-contamination hits and zero final seed-term retention. SRPF yielded the cleanest expository prose and no repetition collapse, but abstracted away from the underwater-machine scene toward generic discourse about intelligent systems, retaining only 0.167 of the seed-term set at the final step.

The results support a restrained conclusion: in this pilot, recursion is sufficient to expose qualitatively distinct attractors in small language models, but none of the tested systems yet justifies strong claims about robust zero-shot knowledge generation. The most promising near-term path is not broader theorizing but tighter experimental control, multi-seed replication, and explicit separation of engine effects from model effects.
\end{abstract}

\section{Introduction}
Large language models are usually evaluated as one-pass next-token predictors. The present project studies a different question: what happens when a model is turned back onto its own outputs through explicit inference-time recursion? Earlier drafts of this project proposed a family of ``hallucination engines'' intended to amplify symbolic structure through repeated self-application. The present paper deliberately narrows the scope. It does not treat prior prose as evidence. It reports only what is directly supported by the executed suite archive supplied for analysis.

The result is therefore a controlled \emph{pilot study}, not a general benchmark. The suite contains four recursion mechanisms paired with three small text-only Hugging Face models, one common seed prompt, one executed run per engine, and structured saved outputs. This is enough to support a careful case study of behavioral differences, failure modes, and immediate methodological lessons. It is \emph{not} enough to support sweeping claims about general model cognition, stable knowledge synthesis, or broad superiority of one recursion family over another.

\section{Related Work}
This pilot sits at the intersection of four strands of prior work. First, it inherits the modern concern with hallucination in large language models, typically framed as a reliability problem rather than a constructive mechanism \cite{ji2023survey,maynez2020faithfulness}. Second, it relates to recursive prompting and self-refinement, where models critique or extend their own prior outputs \cite{wei2022cot,madaan2023selfrefine}. Third, the Activation Gradient Climbing variant is conceptually inspired by DeepDream-style activation maximization \cite{mordvintsev2015inceptionism}. Fourth, the broader framing of recursive self-similarity draws loose inspiration from fractal compression and iterated transform systems \cite{jacquin1992fractal}. 

The contribution here is narrower than a new general theory. It is a cleaned, evidence-backed reconstruction of a speculative project, using real run artifacts to determine which claims survive contact with execution.

\section{Materials and Methods}

\subsection{Verified artifact set}
The analyzed suite archive contains:
\begin{itemize}[leftmargin=1.5em]
    \item a suite manifest recording four executed runs;
    \item a model registry snapshot;
    \item one \texttt{run.json} and one \texttt{run.txt} for each engine;
    \item per-step logged metrics including character count, word count, unique-word ratio, and repeated-bigram count.
\end{itemize}
All results reported below are derived from those saved artifacts.

\subsection{Models and engine assignments}
Table~\ref{tab:design} summarizes the executed design. Two important constraints must be stated immediately. First, only one seed prompt was used. Second, engine and model are partly confounded: RTSM used \texttt{distilgpt2}, AGC used \texttt{gpt2}, and both LVTC and SRPF used \texttt{Qwen2.5-0.5B-Instruct}. Differences therefore reflect end-to-end system behavior, not pure engine effects.

\begin{table}[ht]
\centering
\caption{Executed suite design.}
\label{tab:design}
\begin{tabularx}{\textwidth}{@{}lXlll@{}}
\toprule
Acronym & Engine & Model & Steps & Prompt style \\
\midrule
RTSM & Recursive Token Self-Mirroring & distilgpt2 & 5 & plain \\
AGC & Activation Gradient Climbing & gpt2 & 3 & plain \\
LVTC & Latent Vector Transform Composition & Qwen2.5-0.5B-Instruct & 4 & chat \\
SRPF & Self-Reflective Prompt Feedback & Qwen2.5-0.5B-Instruct & 4 & chat \\
\bottomrule
\end{tabularx}
\end{table}

\subsection{Common seed prompt}
All four systems were initialized from the same seed:
\begin{quote}
\small
\emph{A forgotten machine beneath the sea continued speaking long after its makers were gone. Each repetition made its message stranger, denser, and more self-aware.}
\end{quote}

\subsection{Logged metrics}
The suite already recorded the following per-step metrics in each \texttt{run.json}:
\begin{itemize}[leftmargin=1.5em]
    \item \textbf{character count};
    \item \textbf{word count};
    \item \textbf{type-token ratio} (logged as unique-word ratio);
    \item \textbf{repeated-bigram count}.
\end{itemize}

These are the primary quantitative measures in this paper because they are directly supported by the saved artifacts.

\subsection{Post-hoc exploratory metrics}
To compare the final outputs across engines, three simple post-hoc heuristics were added during analysis:

\paragraph{Seed-term retention.}
From the seed prompt, six content terms were extracted after minimal stop-word removal:
\[
\{\texttt{forgotten}, \texttt{machine}, \texttt{sea}, \texttt{makers}, \texttt{repetition}, \texttt{message}\}.
\]
For each final output, seed-term retention is the fraction of these six terms still present verbatim. This underestimates semantic retention when paraphrase occurs, so it should be read as a conservative lexical proxy.

\paragraph{Adjacent-output Jaccard similarity.}
For each engine, the mean Jaccard similarity between the token sets of consecutive outputs was computed:
\[
J(A,B)=\frac{|A \cap B|}{|A \cup B|}.
\]
High values indicate strong local lexical continuity; low values indicate large step-to-step drift.

\paragraph{Meta-contamination hits.}
A small rule-based detector counted matches to phrases associated with instruction leakage or task-format contamination, such as ``Certainly!'', ``Here's'', ``next recursive transformation'', translation-workflow language, copied control instructions, or explicit assistant-role markers. This heuristic is exploratory only.

\section{Results}

\subsection{Execution summary}
All four runs completed and produced non-empty outputs. The resulting behaviors diverged strongly even under a single shared seed prompt. Table~\ref{tab:summary} gives the main quantitative summary.

\begin{table}[ht]
\centering
\caption{Cross-engine summary from the executed suite.}
\label{tab:summary}
\begin{tabularx}{\textwidth}{@{}lXlrrrrr@{}}
\toprule
Acronym & Engine & Model & Final chars & Final TTR & Min TTR & Max repeated bigrams & Final seed retention \\
\midrule
AGC & Activation Gradient Climbing & gpt2 & 703 & 1.000 & 1.000 & 0 & 0.000 \\
LVTC & Latent Vector Transform Composition & Qwen2.5-0.5B-Instruct & 806 & 0.904 & 0.851 & 2 & 0.000 \\
RTSM & Recursive Token Self-Mirroring & distilgpt2 & 3546 & 0.570 & 0.510 & 236 & 1.000 \\
SRPF & Self-Reflective Prompt Feedback & Qwen2.5-0.5B-Instruct & 691 & 0.885 & 0.878 & 0 & 0.167 \\
\bottomrule
\end{tabularx}
\end{table}

A second view, shown in Table~\ref{tab:drift}, highlights continuity, drift, and contamination.

\begin{table}[ht]
\centering
\caption{Exploratory drift and contamination metrics.}
\label{tab:drift}
\begin{tabularx}{\textwidth}{@{}lXrrr@{}}
\toprule
Acronym & Engine & Mean adjacent Jaccard & Mean TTR & Final meta hits \\
\midrule
AGC & Activation Gradient Climbing & 0.013 & 1.000 & 0 \\
LVTC & Latent Vector Transform Composition & 0.085 & 0.907 & 4 \\
RTSM & Recursive Token Self-Mirroring & 0.778 & 0.804 & 2 \\
SRPF & Self-Reflective Prompt Feedback & 0.085 & 0.927 & 0 \\
\bottomrule
\end{tabularx}
\end{table}

\subsection{Per-engine behavior}

\subsubsection{Recursive Token Self-Mirroring (RTSM)}
RTSM showed the strongest lexical continuity with the seed. Its mean adjacent-output Jaccard similarity was 0.778, and all six seed content terms remained present in the final output. Up through step 3, its type-token ratio stayed high (0.974--0.983) while output length increased steadily. At step 4, however, the run underwent a sharp collapse: character count jumped from 1439 to 3032, type-token ratio dropped to 0.510, and repeated bigrams rose to 235. By the final step, the output had begun copying control text from the experiment itself (including the prompt scaffold ``You are a text-only language model participating in a controlled recursion experiment'').

This is the clearest example in the suite of \emph{motif retention with degenerative recursion}. The system preserved the seed scene, but the recursion loop eventually folded the control prompt back into the generated object-language and entered high-repetition territory.

\subsubsection{Activation Gradient Climbing (AGC)}
AGC behaved very differently. It preserved maximal surface diversity under the logged type-token metric (1.000 at every step) and showed no repeated-bigram buildup, but it also drifted away from the seed almost immediately. The final output contains none of the six seed terms, and the mean adjacent-output Jaccard similarity is only 0.013. Qualitatively, the run moved into fragmented prose involving voices, press-conference narration, and laptop practice sessions, with no stable relation to the underwater machine premise.

In this pilot, AGC looks less like a controlled semantic ascent and more like an unstable attractor-hopping process. It avoids lexical collapse, but it does so by giving up thematic continuity.

\subsubsection{Latent Vector Transform Composition (LVTC)}
LVTC produced moderately diverse prose, with a final type-token ratio of 0.904 and only minor repeated-bigram counts. Its weakness was not repetition but \emph{meta-contamination}. The first output began with an assistant-style compliance marker (``Certainly! Here's the next recursive transformation\ldots''), and later steps drifted further into translation-workflow discourse. The final output discusses ``identifying common topics or themes across sentences'' and ``translating'' segments, retaining none of the six seed content terms. The final output triggered four meta-contamination hits under the exploratory heuristic.

This suggests that, at least in the executed setup, the latent-direction machinery interacted poorly with the instruction-tuned model. Rather than remaining in the fictional/object-language space of the seed, the run repeatedly surfaced into the model's chat-task prior.

\subsubsection{Self-Reflective Prompt Feedback (SRPF)}
SRPF yielded the most readable and compositionally stable prose in the suite. It showed no repeated-bigram collapse, no meta-contamination hits in the final output, and a mean type-token ratio of 0.927. The final text is coherent expository prose about intelligent systems, learning, self-improvement, and technological development.

Its cost was thematic abstraction. Only one of the six seed terms remained detectable at the end, giving a final seed-term retention of 0.167. In other words, SRPF did not preserve the underwater-machine scene as a scene; instead it generalized from the seed into a higher-level discourse about complex learning systems. For some creative or essayistic use-cases this may be desirable. For controlled continuation of a specific motif, it is a drift failure.

\subsection{Per-step trajectories}
Table~\ref{tab:stepmetrics} records the per-step metrics exactly as recovered from the suite logs.

\begin{longtable}{@{}llrrrr@{}}
\caption{Per-step metrics from the saved run logs.}\label{tab:stepmetrics}\\
\toprule
Acronym & Step & Chars & Words & TTR & Repeated bigrams \\
\midrule
\endfirsthead
\toprule
Acronym & Step & Chars & Words & TTR & Repeated bigrams \\
\midrule
\endhead
AGC & 1 & 199 & 39 & 1.000 & 0 \\
AGC & 2 & 634 & 109 & 1.000 & 0 \\
AGC & 3 & 703 & 105 & 1.000 & 0 \\
LVTC & 1 & 652 & 94 & 0.851 & 2 \\
LVTC & 2 & 140 & 18 & 1.000 & 0 \\
LVTC & 3 & 751 & 109 & 0.872 & 2 \\
LVTC & 4 & 806 & 104 & 0.904 & 1 \\
RTSM & 1 & 688 & 116 & 0.974 & 0 \\
RTSM & 2 & 1027 & 171 & 0.982 & 0 \\
RTSM & 3 & 1439 & 236 & 0.983 & 0 \\
RTSM & 4 & 3032 & 494 & 0.510 & 235 \\
RTSM & 5 & 3546 & 570 & 0.570 & 236 \\
SRPF & 1 & 56 & 8 & 1.000 & 0 \\
SRPF & 2 & 563 & 73 & 0.945 & 0 \\
SRPF & 3 & 1010 & 139 & 0.878 & 0 \\
SRPF & 4 & 691 & 87 & 0.885 & 0 \\
\bottomrule
\end{longtable}

\section{Discussion}
Three conclusions survive this pilot.

\paragraph{First, recursion clearly changes model behavior.}
Even with a single seed and small models, the four systems entered visibly different attractors. That alone justifies continuing the project under stricter controls.

\paragraph{Second, preservation and readability pulled in opposite directions.}
RTSM preserved the original motif best, but became the least stable. SRPF produced the cleanest prose, but abstracted away from the scene. In this suite, no engine gave both strong motif retention and strong compositional stability.

\paragraph{Third, some failures came from engine-model interaction rather than recursion alone.}
LVTC's meta-task contamination and AGC's thematic drift cannot be cleanly interpreted as engine-only failures because the engines were not run across a balanced model grid. This is especially important for instruction-tuned models, whose built-in chat priors may dominate fragile recursion schemes.

\section{Limitations}
This study has substantial limitations, and they should be treated as design constraints for the next round rather than as footnotes.

\begin{enumerate}[leftmargin=1.5em]
    \item \textbf{Single seed prompt.} The entire suite is based on one fictional prompt. No claims about average behavior across tasks are justified.
    \item \textbf{One run per engine.} There is no statistical replication across random seeds.
    \item \textbf{Engine-model confound.} Most engines were paired with different models, so engine effects and model effects are not separable.
    \item \textbf{Small-model regime only.} All tested models were relatively small and CPU-run.
    \item \textbf{Exploratory post-hoc metrics.} Seed-term retention, adjacent Jaccard, and meta-contamination were added during analysis and should be treated as heuristic rather than canonical measures.
    \item \textbf{Lexical rather than semantic continuity.} The current analysis does not include embedding-based semantic scoring or human evaluation.
\end{enumerate}

\section{Conclusion}
The executed suite provides a real foothold for the project, but it narrows the defensible claim substantially. On the evidence available here, latent recursion in small language models does \emph{not} yet look like a robust zero-shot knowledge generator. What it \emph{does} look like is a family of controllable failure modes and attractor behaviors: self-mirroring with prompt-copy collapse, activation-based drift, latent-direction meta-task contamination, and reflective feedback that improves prose while generalizing away from the source scene.

That is still a meaningful result. It shows that the project can move forward honestly, using real artifacts rather than speculative prose. The next step is straightforward: rerun the suite on a balanced engine--model grid, with multiple seeds, fixed random seeds for replication, and an evaluation stack that combines lexical metrics, semantic similarity, and blinded human judgments of coherence, originality, and motif retention.

\appendix
\section{Short qualitative excerpts}
The following short excerpts are included only to anchor the quantitative interpretation.

\begin{itemize}[leftmargin=1.5em]
    \item \textbf{AGC final output:} ``that time because it was late starting out \ldots on their laptop \ldots''
    \item \textbf{LVTC final output:} ``The next recursive transformation involves identifying common topics or themes \ldots before translating.''
    \item \textbf{RTSM final output:} ``You are a text-only language model participating in a controlled recursion experiment.''
    \item \textbf{SRPF final output:} ``the recognition of these systems' capacity for learning and self-improvement \ldots''
\end{itemize}

\section{Reproducibility note}
This paper was written from the saved suite archive only. It intentionally avoids treating earlier draft prose as experimental evidence. Any future extension of the paper should preserve that rule.

\begin{thebibliography}{99}

\bibitem{brown2020gpt3}
Tom B. Brown, Benjamin Mann, Nick Ryder, Melanie Subbiah, Jared Kaplan, Prafulla Dhariwal, Arvind Neelakantan, et al.
\newblock Language Models are Few-Shot Learners.
\newblock \emph{Advances in Neural Information Processing Systems}, 33, 2020.

\bibitem{ji2023survey}
Zhemin Ji, Nayeon Lee, Rita Frieske, Tiantian Yu, Dan Su, Yan Xu, Etsuko Ishii, et al.
\newblock Survey of Hallucination in Natural Language Generation.
\newblock \emph{ACM Computing Surveys}, 55(12), 2023.

\bibitem{jacquin1992fractal}
Arnaud E. Jacquin.
\newblock Image Coding Based on a Fractal Theory of Iterated Contractive Image Transformations.
\newblock \emph{IEEE Transactions on Image Processing}, 1(1):18--30, 1992.

\bibitem{madaan2023selfrefine}
Aman Madaan, Niket Tandon, Prakhar Gupta, Peter Clark, Yiming Yang, Shashank Srivastava, and others.
\newblock Self-Refine: Iterative Refinement with Self-Feedback.
\newblock \emph{arXiv preprint arXiv:2303.17651}, 2023.

\bibitem{maynez2020faithfulness}
Joshua Maynez, Shashi Narayan, Bernd Bohnet, and Ryan McDonald.
\newblock On Faithfulness and Factuality in Abstractive Summarization.
\newblock \emph{Proceedings of ACL}, 2020.

\bibitem{mikolov2013efficient}
Tomas Mikolov, Ilya Sutskever, Kai Chen, Greg Corrado, and Jeffrey Dean.
\newblock Distributed Representations of Words and Phrases and their Compositionality.
\newblock \emph{Advances in Neural Information Processing Systems}, 2013.

\bibitem{mordvintsev2015inceptionism}
Alexander Mordvintsev, Christopher Olah, and Mike Tyka.
\newblock Inceptionism: Going Deeper into Neural Networks.
\newblock Google Research Blog, 2015.

\bibitem{wei2022cot}
Jason Wei, Xuezhi Wang, Dale Schuurmans, Maarten Bosma, Brian Ichter, Fei Xia, Ed Chi, Quoc Le, and Denny Zhou.
\newblock Chain-of-Thought Prompting Elicits Reasoning in Large Language Models.
\newblock \emph{Advances in Neural Information Processing Systems}, 2022.

\end{thebibliography}

\end{document}
